\begin{center}\begin{tikzpicture}[x=1.45cm,y=1.45cm,>=stealth,font=\small]
\node[font=\small\bfseries] at (0,3.15) {中点确定弦方向：$AM=MB$};
\draw[->,gray] (-3.5,0)--(3.7,0) node[right] {$x$};
\draw[->,gray] (0,-2.4)--(0,2.5) node[above] {$y$};
\draw[CStatement,thick] (0,0) ellipse (3 and 2);
\coordinate (A) at (3,0); \coordinate (B) at (0,2); \coordinate (M) at (1.5,1);
\draw[teal,very thick] (A)--(B);
\fill (A) circle(1.8pt) node[below right] {$A$};
\fill (B) circle(1.8pt) node[above left] {$B$};
\fill[orange!80!black] (M) circle(2pt) node[above right] {$M$};
\node[below left] at (0,0) {$O$};
\draw[orange!80!black,dashed,thick] (0,0)--(M);
\node[below,align=center] at (0,-2.85) {$x^2/9+y^2/4=1,\quad M(3/2,1),\quad k_{AB}=-2/3$};

\end{tikzpicture}
\end{center}
\begin{statementbox}
\textbf{\textcolor{CStatement}{条件}}
\begin{description}[style=nextline,leftmargin=2.8em,labelsep=0.5em,itemsep=0.35em]
  \item[\textbf{条件 1（椭圆）：}] $\dfrac{x^2}{a^2}+\dfrac{y^2}{b^2}=1$，$a>b>0$。
  \item[\textbf{条件 2（真实弦）：}] $A,B$ 是椭圆上两个不同点，$M(x_0,y_0)$ 为弦 $AB$ 的中点；于是 $\dfrac{x_0^2}{a^2}+\dfrac{y_0^2}{b^2}<1$。
  \item[\textbf{条件 3（求斜率时）：}] $y_0\ne0$。此时弦必不竖直，斜率存在。
\end{description}
\medskip
\textbf{\textcolor{CStatement}{结论一（中点弦斜率）}}
\[
\boxed{\quad k_{AB}=-\frac{b^2}{a^2}\frac{x_0}{y_0}
             =-\frac{b^2x_0}{a^2y_0}\quad(y_0\ne0).\quad}
\]
\textbf{\textcolor{CStatement}{结论二（更完整：不除以 $x_0,y_0$ 的弦方程）}}
若 $M\ne O(0,0)$，则弦 $AB$ 所在直线为
\[
\boxed{\quad\frac{x_0}{a^2}(x-x_0)+\frac{y_0}{b^2}(y-y_0)=0.\quad}
\]
这个方程既能处理一般斜弦，也能处理水平弦、竖直弦。对椭圆内部的每个非中心点 $M$，此式都确定\textbf{唯一}的一条以 $M$ 为中点的非退化弦。
\medskip
\textbf{\textcolor{CStatement}{特殊位置（不能机械套斜率式）}}
\begin{description}[style=nextline,leftmargin=2.8em,labelsep=0.5em,itemsep=0.3em]
\item[\textbf{$x_0=0,\ y_0\ne0$：}] $k=0$，弦水平，方程 $y=y_0$。
\item[\textbf{$y_0=0,\ x_0\ne0$：}] 弦竖直，方程 $x=x_0$，斜率不存在。
\item[\textbf{$M=O$：}] 上述无除法方程退化为 $0=0$；任意过中心的直线均可截得中心弦，无法由中点唯一确定斜率。
\end{description}
\medskip
\textbf{\textcolor{CStatement}{常见变形（焦点在 $y$ 轴）}}
对于 $\dfrac{x^2}{b^2}+\dfrac{y^2}{a^2}=1$（$a>b>0$），
\[
\boxed{\quad k=-\frac{a^2x_0}{b^2y_0}\quad(y_0\ne0).\quad}
\]
\textbf{统一辨认法：}把两个分母直接记成 $A,B>0$，即 $\frac{x^2}{A}+\frac{y^2}{B}=1$，则
\[
\boxed{\quad\frac{x_0}{A}(x-x_0)+\frac{y_0}{B}(y-y_0)=0,
\qquad k=-\frac{Bx_0}{Ay_0}.\quad}
\]
避免只按“长轴在 $x$ 轴还是 $y$ 轴”死记两套公式。
\end{statementbox}
